\begin{tabularx}{\textwidth}{L{31mm}X X}
\toprule
\textbf{Fogalom} & \textbf{Jelentése} & \textbf{Vizsgán fontos megjegyzés} \\
\midrule
Osztály & Programozó által definiált típus, amely adatokból és műveletekből épül fel. & Az osztály saját névteret és hatáskört képez; Java-ban a program lényegében osztálydefiníciók összessége. \\
Objektum & Egy osztály konkrét példánya. & Több objektum készülhet ugyanabból az osztályból: a viselkedésük azonos jellegű, de az állapotuk eltérhet. \\
Adattag & Az objektum vagy az osztály állapotát tároló változó. & Az információrejtés miatt az állapotot általában \code{private} mezőkben tároljuk. \\
Metódustag & Az osztályhoz tartozó művelet. & A publikus metódusok adják az objektum külső felületét, vagyis azt, ahogyan más objektumok üzenetet küldenek neki. \\
Konstruktor & Objektum létrehozásakor lefutó inicializáló tag. & Nem visszatérési típusos metódus; feladata, hogy az objektum érvényes kezdőállapotba kerüljön. \\
Referencia & Objektumra mutató érték, amelyen keresztül az objektum elérhető. & A referencia másolása nem objektummásolás: több referencia is ugyanarra az objektumra mutathat. \\
\bottomrule
\end{tabularx}
